\section{Agentic Retrieval and Generation}
\label{sec:online}

In the agentic retrieval and generation stage, our primary objectives are to effectively retrieve relevant information using the constructed EKG and to generate robust, contextually accurate responses based on the retrieved data. To achieve this, we introduce an agentic searching mechanism that explores multiple retrieval pathways within the index to gather the necessary information. Additionally, we propose a thought consistency strategy, enabling \sysname to generate and evaluate multiple responses from different retrieval pathways, ultimately selecting the most appropriate one.

\subsection{Tri-View Retrieval}
\label{subsec:tri_view}

\begin{figure}
    \centering
    \includegraphics[width=0.9\linewidth]{Figures/borda_counting-crop.pdf}
    \caption{An illustration of tri-view retrieval and borda counting on the retrieved events.}
    \label{fig:borda_counting}
\end{figure}


To comprehensively retrieve relevant information from the index for a given query, \sysname employs a tri-view retrieval process: the first view targets events, enabling retrieval at the event level to provide information for event summary-related queries. Specifically, the query is encoded using the text encoder JinaCLIP~\cite{koukounas2024jinaclip} and matched against the events table in the constructed EKG. The second view focuses on entities, offering insights into basic facts or item-specific queries. For this, we leverage the entity centroids extracted and aggregated as detailed in \S\ref{subsec:entity_extraction} to facilitate retrieval. The third view utilizes vision embeddings of raw video frames as complementary information. The retrieved entities and raw frames are subsequently linked to their associated events through the constructed EKG.







It is important to note that the retrieved events should be ranked. Ranking is not only crucial for filtering noise from the retrieved results but also essential for enabling agentic searching, as detailed in \S\ref{subsec:agentc_searching}. A straightforward ranking method, such as similarity-based ranking, cannot be directly applied to \sysname due to the retrieved events originating from three distinct views. To integrate these results, we propose to use a weighted Borda counting approach.

Fig.~\ref{fig:borda_counting} illustrates the process of using Borda counting to integrate and rank retrieved events from the three views in \sysname. Specifically, we select the top $K$ events from each view and rank them based on their calculated similarities within that view. Subsequently, the similarities of these $K$ events are normalized to compute their Borda scores:

\begin{equation}
    s_{m}(e_j) = \frac{\text{sim}_m(e_j)}{\sum_{e_k \in \mathcal{E}_m} \text{sim}_m(e_k)},
\end{equation}
where $\mathcal{E}_m$ represents the set of events retrieved from view $m$. The final Borda score for each event $e_j$ is then obtained by summing its scores across all views:

\begin{equation}
    s(e_j) = \sum_{m} s_m(e_j),
\end{equation}

Finally, the aggregated Borda scores $s(e_j)$ are used to rank all retrieved events.






\begin{figure}
    \centering
    \includegraphics[width=0.9\linewidth]{Figures/tree_search-crop.pdf}
    \caption{An example of agentic tree search with four actions and a depth of three, yielding 13 distinct pathways for information gathering and response generation.}
    \label{fig:tree_search}
\end{figure}

\subsection{Agentic Searching on Graph}
\label{subsec:agentc_searching}




The retrieved events mentioned above can be directly utilized to generate the final answer. However, to support complex queries, such as query-focused summaries and multi-hop queries, \sysname searches for additional relevant information by leveraging the relationships between events and entities within the constructed EKG. To enable efficient exploration, we propose the agentic searching on graph approach.

\revised{
The proposed approach is inspired from human strategies for information retrieval and reasoning within video content. Typically, individuals initiate this process by locating salient segments guided by retrieval keywords (\ie query), subsequently enriching their understanding by examining the temporal context in adjacent clips (\ie forward and backward). For a more thorough comprehension, they may iteratively refine their search using alternative keywords as necessary (\ie re-query). Similarly, in our agentic searching process, we define the agentic action space as follows:
}



    
\textbf{Forward (F)}: this action extends the current retrieval by including temporally subsequent events on the EKG for all events in the event list. It reflects the natural tendency of humans to seek forward narrative progression when trying to understand what happens next or how a situation evolves over time.
    
\textbf{Backward (B)}: complementing the forward action, this action retrieves temporally preceding events, enabling a backward exploration of the narrative to uncover prior context or causal factors.
    
\textbf{Re-query (RQ)}: this action generates a new query represented by a list of keywords via an LLM and retrieves complementary events as outlined in \S\ref{subsec:tri_view}. It reflects the human tendency to gather information from multiple perspectives to achieve a more comprehensive understanding.

\textbf{Summary and Answer (SA)}: this action utilizes the descriptions of the retrieved events from the EKG and generates the response to the specific query by employing an LLM.







\revised{
Utilizing these predefined agentic actions, we structure the agentic search process within a tree search framework. The search is initiated by performing an initial retrieval based on the original query, yielding a set of relevant events that constitute the root node of the search tree, as outlined in \S\ref{subsec:tri_view}. For each node, a single-step rollout is conducted in which four predefined actions, namely, \textit{forward}, \textit{backward}, \textit{re-query}, and \textit{summary and answer}, are executed on every events on the node. Upon reaching the {SA} action, the corresponding search trajectory is terminated by generating an answer. The rollout proceeds iteratively until the predefined maximum tree depth is reached. This tree search mechanism systematically explores multiple pathways to extract information from the EKG. As depicted in Fig.~\ref{fig:tree_search}, a tree of depth 3 yields 13 distinct information-gathering paths along with their respective answers. All generated answers are subsequently evaluated, with the optimal answer selected via the thought consistency method elaborated in \S\ref{subsec:thoughts_consistency}.
}




A practical issue in the tree search process is the exponential growth in the number of retrieved events as the tree depth increases. This not only introduces computational overhead but also results in the accumulation of noisy or irrelevant information. To mitigate this, we use a length constraint on the maintained event list during the search process, \ie 16 in our implementation. When the number of retrieved events exceeds this limit, we employ a drop strategy to discard less relevant events based on their rankings described in \S\ref{subsec:tri_view}.




\subsection{Consistency Enhanced Generation}
\label{subsec:thoughts_consistency}


During the agentic tree search, multiple candidate answers are generated at {SA} nodes across different pathways. To determine the final answer, it is necessary to either select or synthesize from these candidates. A straightforward approach would be majority voting. However, due to the diversity of retrieval paths, only a small subset of these nodes is likely to access the essential information with minimal noise, producing high-quality answers. To this end, we introduce the thoughts-consistency mechanism to identify and select the most reliable final answer.




At each SA node, instead of generating the answer a single time, we repeatedly generate answers multiple times using a Chain-of-Thought (CoT) prompting scheme. Following the principle of self-consistency~\cite{wang2023self}, correct answers are more likely to emerge consistently across multiple valid reasoning trajectories during repeated generations. Specifically, we evaluate the consistency not only across the generated answers but also within their associated CoT traces. Nodes demonstrating strong internal coherence, where the reasoning process aligns logically with the conclusion, are assigned higher scores.


To formalize this process, we propose a scoring framework that integrates both \emph{answer agreement} and \emph{thought consistency}. At each SA node, we perform $n$ rounds of sampling using a temperature setting between $0.5$ and $0.7$, resulting in a set of $n$ candidate outputs denoted as $\{(a_i, r_i)\}_{i=1}^n$, where $a_i$ is the answer and $r_i$ is the associated reasoning trace.

Let $\mathcal{A} = \{a^{(1)}, a^{(2)}, \dots, a^{(T)}\}$ be the set of unique answers among the $n$ samples, where $T$ is the number of distinct answers. The answer agreement score $S_a^{(t)}$ for a candidate answer $a^{(t)}$ is defined as the proportion of times it appears in the samples:
\begin{equation}
S_a^{(t)} = \frac{|\{i \mid a_i = a^{(t)}\}|}{n}
\end{equation}

The thought consistency score $S_r^{(t)}$ for $a^{(t)}$ is computed as the average BERTScore between all pairs of reasoning traces associated with $a^{(t)}$:
\begin{equation}
S_r^{(t)} = \frac{2}{k(k-1)} \sum\limits_{1 \leq i < j \leq k} \text{BERTScore}(r_i, r_j),
\end{equation}
where $k$ is the number of times $a^{(t)}$ appears in the $n$ samples. The final score for each candidate answer combines these two components:
\begin{equation}
\label{eq:thoughts_consistency_weight}
S_{\text{final}}^{(t)} = \lambda S_a^{(t)} + (1 - \lambda) S_r^{(t)},
\end{equation}
where $\lambda \in [0, 1]$ is a weighting parameter controlling the trade-off between answer agreement and thought consistency. In our implementation, we set it to 0.3, the parameter tuning would be discussed in \S\ref{sec:thoughts_consistency}.



For each {SA} node, the candidate answer with the highest $S_{\text{final}}^{(t)}$ is selected as its definitive response. To enhance the reliability of this final answer, we propose an additional agentic action, \textbf{Check Frames and Answer (CA)}. This action retrieves the raw video frames associated with the events from the EKG and utilizes the VLM to generate a refined response to the specific query. By doing so, this action effectively supplements any missing information relevant to the query that may have been overlooked during the construction phase.


Specifically, after ranking all candidate answers from the {SA} nodes using the consistency-enhanced scoring mechanism, the top-2 nodes with differing answers are selected. The video frames corresponding to their retrieved events are extracted, and the VLM is prompted to generate a new response by directly attending to the visual evidence. Furthermore, the thought-consistency mechanism is applied to the {CA} nodes to bolster the reliability of the final generated answer.


